\chapter{Active Circuits, Feedback, and Oscillators}
\label{ch:active}
\section{Why Active Circuits Appear}
The Extra exam includes amplifiers, op-amps, feedback, oscillators, mixers, and
regulated power supplies. A full electronics text is not needed here; the goal is
to connect the circuit formulas to system behavior---and, in this chapter, to the
classical feedback ideas that a controls background already contains. The device
details (bias points, transistor models, tube parameters) live in the ARRL
manual; what follows is the systems view that ties them together.

\section{Amplifier Gain and Classes}
Voltage gain is \(A_v=V_o/V_i\); in decibels \(A_{v,\dB}=20\log_{10}|A_v|\). Power
gain is \(G_{\dB}=10\log_{10}(P_o/P_i)\).

Amplifier \emph{class} describes how much of the input cycle the active device
conducts, which trades linearity against efficiency:
\begin{description}
\item[Class A] conducts the full \(360^\circ\); most linear, least efficient
(theoretical max \SI{50}{\percent} transformer-coupled, near \SI{25}{\percent}
resistively coupled).
\item[Class B / AB] Class B conducts half the cycle (\(180^\circ\)); Class AB
somewhat more than half. Higher efficiency than Class A, with some crossover
distortion; common in linear SSB power amplifiers.
\item[Class C] conducts less than half the cycle; efficient but nonlinear, used
where a tuned output tank restores the sinusoid (FM, CW).
\end{description}
Class C is nonlinear on purpose: the output tank is a high-\(Q\) resonator (a
pole pair near the \(\jj\omega\) axis) that rings at the fundamental and rejects
the harmonics the clipped current injects. The same nonlinearity that spoils
linearity is what makes mixers and frequency multipliers work.

\section{Op-Amp Ideals}
An ideal op-amp has very high input impedance, very low output impedance, and
very large open-loop gain \(A(s)\). Negative feedback makes practical closed-loop
gains predictable:
\[
\text{inverting:}\quad \frac{V_o}{V_i}=-\frac{R_f}{R_{in}},
\qquad
\text{non-inverting:}\quad \frac{V_o}{V_i}=1+\frac{R_f}{R_g}.
\]
These come from the two ideal rules---no current into the inputs, and the inputs
driven to equal voltage by feedback---but both rules are consequences of a deeper
statement: the loop gain is large.

\section{Feedback as a Loop}
Model the amplifier as a forward path \(A(s)\) with a fraction \(\beta\) of the
output fed back and subtracted at the input.

\begin{figure}[htbp]
\centering
\begin{tikzpicture}[>=Latex,node distance=1.1cm,scale=1]
\node (in) at (0,0) {\(V_i\)};
\node[draw,circle,inner sep=1.4pt] (sum) at (1.3,0) {\(\Sigma\)};
\node[draw,minimum width=1.3cm,minimum height=0.8cm] (A) at (3.6,0) {\(A(s)\)};
\node (out) at (6.3,0) {\(V_o\)};
\node[draw,minimum width=1.1cm,minimum height=0.7cm] (B) at (3.6,-1.5) {\(\beta\)};
\draw[->] (in) -- (sum);
\draw[->] (sum) -- (A);
\draw[->] (A) -- (out);
\draw[->] (5.3,0) |- (B);
\draw[->] (B) -| (sum);
\node at (1.15,0.32) {\small \(-\)};
\end{tikzpicture}
\caption{Standard negative-feedback loop. The loop gain is \(L(s)=A(s)\beta\).}
\label{fig:feedback-loop}
\end{figure}

The closed-loop gain is
\[
A_{\text{cl}}(s)=\frac{A(s)}{1+A(s)\beta}
=\frac{A(s)}{1+L(s)},
\qquad
L(s)=A(s)\beta .
\]
When the loop gain \(|L|\gg1\), \(A_{\text{cl}}\to 1/\beta\): the gain is set by
the passive feedback network, not by the messy device parameters. That single
fact is why the inverting and non-inverting formulas above depend only on
resistor ratios.

\begin{controlsbox}
\textbf{Refresher: loop gain and the sensitivity trade.} \(L(s)=A(s)\beta\) is the
gain around the loop with it broken open. Everything classical control says about
a unity-feedback loop applies. Large \(|L|\) desensitizes the closed-loop gain to
changes in \(A\): a fractional change \(\dd A/A\) becomes \(\dd A_{\text{cl}}/
A_{\text{cl}}=(\dd A/A)/(1+L)\). You are spending loop gain to buy immunity from
device variation, temperature, and distortion---the same currency you spend in a
process-control loop.
\end{controlsbox}

\section{Gain--Bandwidth Product}
Model the op-amp open-loop response as a single dominant pole,
\[
A(s)=\frac{A_0}{1+s/\omega_a},
\]
with large DC gain \(A_0\) and a low open-loop corner \(\omega_a\). Close the loop
with constant \(\beta\):
\[
A_{\text{cl}}(s)=\frac{A_0}{1+A_0\beta}\cdot
\frac{1}{1+s/\bigl[\omega_a(1+A_0\beta)\bigr]}.
\]
The closed-loop DC gain drops by \((1+A_0\beta)\) and the closed-loop corner rises
by the same factor. Their product is invariant:
\[
(\text{gain})\times(\text{bandwidth})=A_0\,\omega_a=\text{constant}.
\]
This is the \emph{gain--bandwidth product}. In pole language, negative feedback
slides the single pole to the left---faster response, wider bandwidth---at the
cost of gain. Doubling closed-loop gain halves closed-loop bandwidth.

\begin{controlsbox}
\textbf{Refresher: pole motion under feedback.} Wrapping feedback around
\(A_0/(1+s/\omega_a)\) puts the closed-loop pole at
\(s=-\omega_a(1+A_0\beta)\). As you increase \(\beta\), the pole travels left
along the real axis---a one-branch root locus. A single pole can never be
destabilized by more feedback; it just gets faster. Instability needs at least
two poles, which is the next section.
\end{controlsbox}

\section{Stability, Phase Margin, and Gain Margin}
Real amplifiers have more than one pole. With two or more poles the loop gain
\(L(\jj\omega)\) accumulates phase lag, and if the phase reaches \(-180^\circ\)
while \(|L|\ge 1\), the subtracting junction effectively adds---the feedback has
become positive at that frequency and the closed-loop poles move toward, or
into, the right half-plane.

\begin{controlsbox}
\textbf{Refresher: gain and phase margin.} Evaluate the loop gain
\(L(\jj\omega)\) around the loop.
\begin{itemize}
\item The \emph{gain-crossover} frequency is where \(|L(\jj\omega)|=1\)
(\(\SI{0}{\decibel}\)). The \emph{phase margin} is how far the phase there is
above \(-180^\circ\).
\item The \emph{phase-crossover} frequency is where \(\angle L=-180^\circ\). The
\emph{gain margin} is how far \(|L|\) there is below \(\SI{0}{\decibel}\).
\end{itemize}
Positive margins mean the closed-loop poles stay in the left half-plane. Small
margins mean poles near the \(\jj\omega\) axis: ringing, overshoot, a peaked
frequency response. Zero margin means poles on the axis---the boundary of
oscillation. This is the Nyquist/Bode stability story from a classical controls
course, applied to an RF or audio amplifier instead of a chemical process.
\end{controlsbox}

Compensation (deliberately lowering the first pole, as in an internally
compensated op-amp) restores phase margin by making the loop gain fall through
\SI{0}{\decibel} before the second pole adds its lag. It buys stability with
bandwidth---the same trade, again.

\section{Oscillators: Poles Placed on the Axis}
An oscillator is a feedback loop designed to sit exactly on the stability
boundary. The closed-loop poles are the roots of \(1+L(s)=0\), so they reach the
\(\jj\omega\) axis when
\[
L(\jj\omega_{\text{osc}})=-1,
\qquad\text{that is}\qquad
|L(\jj\omega_{\text{osc}})|=1,\quad
\angle L(\jj\omega_{\text{osc}})=-180^\circ,
\]
the Nyquist \(-1\) point. This is the \emph{Barkhausen criterion}. Its familiar
``unit magnitude, zero net phase'' wording refers to the round-trip transmission
\(T=-L\), which already includes the \(-1\) of the subtracting junction; setting
\(T=1\) is the same statement as \(L=-1\). In s-plane language the closed-loop
poles are placed on the \(\jj\omega\) axis at \(\pm\jj\omega_{\text{osc}}\)---the
marginally stable case from the s-plane chapter. At start-up the design makes \(|L|>1\) so the pole pair sits slightly in
the right half-plane and the amplitude grows; a nonlinearity (device saturation,
a limiter, a lamp or AGC) then reduces the effective gain until \(|L|=1\) and the
poles settle onto the axis at a stable amplitude.

The frequency-setting network sets where those poles land:
\begin{description}
\item[LC oscillators] (Colpitts, Hartley) use an LC tank; \(\omega_{\text{osc}}
\approx 1/\sqrt{LC}\). Higher tank \(Q\) means poles closer to the ideal axis
location and lower phase noise.
\item[Crystal oscillators] replace the tank with a quartz resonator of very high
\(Q\), giving excellent frequency stability.
\item[RC phase-shift / Wien-bridge] use RC networks to reach the required loop
phase; common at audio frequencies.
\end{description}

\begin{physicalbox}
Phase noise is the frequency-domain fingerprint of how firmly the poles are
pinned to the axis. A high-\(Q\) resonator holds the pole pair tightly at
\(\pm\jj\omega_{\text{osc}}\), so the spectrum is a sharp line; a low-\(Q\) loop
lets the poles wander, smearing the line. This is why crystal and high-\(Q\) LC
oscillators are quieter.
\end{physicalbox}

\section{Active Filters and Their Limits}
Op-amps can implement active RC filters that synthesize the second-order pole
pairs of \cref{ch:filters} without inductors, by using feedback to create complex
poles from resistors and capacitors (Sallen--Key, multiple-feedback). The catch
is that the same finite gain--bandwidth product limits them: near the amplifier's
own gain-crossover the ideal-op-amp assumption fails, so an active filter that
behaves at audio may be useless at RF, where LC or crystal filters take over.

\section{Power Supplies}
Rectifiers, filter capacitors, bleeder resistors, regulators, and switching
converters reuse the elements from earlier chapters. A linear regulator is itself
a feedback loop---it compares the output to a reference and drives the pass
element to null the error, with the same stability considerations as any feedback
system. Equal-value resistors across series filter capacitors equalize the shared
voltage and discharge stored energy after power is removed.

\section{Exam Relevance}
\begin{exambox}
For active-circuit questions, keep three layers separate: small-signal gain,
power handling, and stability. Feedback questions almost always reduce to one
idea---\emph{loop gain}: large loop gain buys predictable gain, wide bandwidth,
and low distortion; too much phase lag at unity loop gain costs stability; and an
oscillator is the deliberate limiting case with the poles parked on the
\(\jj\omega\) axis.
\end{exambox}
